\documentclass[a4paper,11pt]{article}
\usepackage{jcappub} % for details on the use of the package, please see the JINST-author-manual
\usepackage{lineno}
\usepackage{pifont}% http://ctan.org/pkg/pifont
\usepackage{array}
\usepackage{siunitx}
\usepackage{multirow}
\usepackage{graphicx}
%\usepackage{floatrow}
\usepackage{subfig}
\usepackage{multicol}
%\usepackage{subcaption}
\usepackage{comment}
\usepackage{dsfont}
\usepackage{soul}
%%%%%%%%%%%%%%%%%%%%%%%%%%%%%%%%%%%%%%%%
\usepackage{txfonts}
\usepackage{color}
\usepackage[numbers]{natbib}

%\linenumbers
\newcommand{\xmark}{\ding{55}}
\newcommand{\cmark}{\ding{51}}
\newcommand{\todo}[1]{{\textcolor[rgb]{1, 0, 1}{[TODO: #1]}}}

%\arxivnumber{arXiv:XXXX.XXXXX} % Only if you have one
\title{\boldmath Spectral Imaging with QUBIC: Atmospheric Reconstruction from Time-Ordered-Data using Bolometric Interferometry}

% Collaborations

%% [A] If main author
%% \collaboration{\includegraphics[height=17mm]{collabroation-logo}\\[6pt]
%%  XXX collaboration}

%% or
%% [B] If "on behalf of"
%% \collaboration[c]{on behalf of XXX collaboration}


% Authors
% The "\note" macro will give a warning: "Ignoring empty anchor...", you can safely ignore it.

%% [A] simple case: 2 authors, same institution
%% \author[1]{A. Uthor\note{Corresponding author.}}
%% \author{and A. Nother Author}
%% \affiliation{Institution,\\Address, Country}

%% or, e.g.
%% [B] more complex case: 4 authors, 3 institutions, 2 footnotes
%% \author[a,b]{F. Irst,\}
%% \author[c]{S. Econd,}
%% \author[a,1]{T. Hird\note{Also at Some University.}}
%% \author[c,1]{and Fourth}
%% \affiliation[a]{Institution_1,\\Address, Country}
%% \affiliation[b]{Institution_2,\\Address, Country}
%% \affiliation[c]{Institution_3,\\Address, Country}

\author[a]{T.~Laclavere,}
%% TODO: complete author list and affiliations (see FMM/CMM main.tex for the current QUBIC collaboration list)
\affiliation[a]{Université de Paris, CNRS, Astroparticule et Cosmologie, F-75006 Paris, France}

% E-mail addresses: only for the corresponding author
\emailAdd{laclavere@apc.in2p3.fr}

\abstract{\todo{Context: 1-2 sentences on B-mode searches and why atmospheric contamination from ground-based CMB observations, and QUBIC in particular, matters.}
  % aims heading (mandatory)
   {\todo{Aims: state that this article presents a method to reconstruct and separate the atmospheric signal from the sky signal directly at the Time-Ordered-Data level using Bolometric Interferometry.}}
  % methods heading (mandatory)
   {\todo{Methods: summarize the acquisition model (coordinate system, wind operator), the atmosphere model/simulations, and the reconstruction/separation algorithm.}}
  % conclusions heading (optional), leave it empty if necessary
   {\todo{Conclusions: summarize the main quantitative result(s).}}}


\begin{document}
\maketitle
\flushbottom

\section{Introduction}
\label{sec:introduction}

\subsection{Motivations}
\label{sec:motivations}

\todo{Motivate the need to reconstruct and remove the atmospheric contribution from ground-based CMB observations, and why this is a specific challenge/opportunity for QUBIC. Cite the QUBIC instrument papers~\citep{2020.QUBIC.PAPER1, 2020.QUBIC.PAPER2,2020.QUBIC.PAPER3, 2020.QUBIC.PAPER4, 2020.QUBIC.PAPER5, 2020.QUBIC.PAPER6, 2020.QUBIC.PAPER7, 2020.QUBIC.PAPER8} as well as the companion Frequency Map-Making~\citep{fmm} (hereafter FMM) and Component Map-Making~\citep{cmm} (hereafter CMM) articles.}

\subsection{Goals}
\label{sec:goals}

\todo{State the goals of this article: reconstructing frequency (or component) maps of the atmosphere jointly with, or separately from, the sky signal, and/or reconstructing the wind, using Bolometric Interferometry's spectral-imaging capability.}

\section{Bolometric Interferometry}
\label{sec:bi}

\todo{Short reminder of Bolometric Interferometry, as introduced in FMM section~2 and CMM section~2, adapted to the atmospheric reconstruction context.}

\subsection{Instrument}
\label{sec:instrument}

\todo{Recall the QUBIC instrument main parameters (see FMM Table~1).}

\subsection{Synthesized beam}
\label{sec:synthesized_beam}

\todo{Recall the frequency dependence of the QUBIC synthesized beam (see FMM section~2.2) and why it is relevant to separate the atmosphere from the sky signal.}

\subsection{Map-Making principle (with BI)}
\label{sec:mapmaking_principle}

\todo{Recall the general inverse-problem map-making framework (see FMM section~2.1 and CMM section~2.2) that will be extended in section~\ref{sec:algorithm} to include the atmosphere.}

\section{Atmospheric model and simulations}
\label{sec:atm_model}

\subsection{Atmospheric properties}
\label{sec:atm_properties}

\todo{Describe the physical properties of the atmospheric signal relevant to QUBIC (emission mechanism, frequency dependence, spatial/temporal correlation structure, typical amplitude).}

\subsection{Numerical simulations}
\label{sec:atm_simulations}

\todo{Describe how the atmosphere is numerically simulated (e.g.\ frozen screen advected by wind, power spectrum used, software/tools).}

\subsection{Atmospheric correlations (Imager vs BI)}
\label{sec:atm_correlations}
%% Note: in the outline this subsection was circled with a question mark
%% suggesting it might instead be moved to the appendix -- keeping it here
%% as a regular subsection for now, move to appendix if it ends up being a
%% secondary/technical point rather than a main result.

\todo{Compare the correlation structure of the atmospheric signal as seen by a classical imager versus a Bolometric Interferometer.}

\section{Atmospheric reconstruction and separation algorithm}
\label{sec:algorithm}

\subsection{Data acquisition model}
\label{sec:acquisition_model}

\todo{Extend the acquisition model of Eq.~(FMM~2.1)/(CMM~2.1) to include an atmospheric contribution alongside the sky signal.}

\subsubsection{Coordinate system}
\label{sec:coord_system}

\todo{Describe the coordinate system used to describe the atmospheric screen with respect to the instrument's pointing (e.g.\ horizontal/az-el frame vs.\ sky frame).}

\subsubsection{Wind operator}
\label{sec:wind_operator}

\todo{Describe the operator modeling the advection of the atmospheric screen by the wind across the field of view during the observation.}

\subsection{Atmospheric separation}
\label{sec:atm_separation}

\todo{Describe how the atmospheric component is separated from the sky signal within the map-making/component-separation framework (see CMM section~5 for the analogous astrophysical component separation algorithm).}

\subsection{Wind reconstruction}
\label{sec:wind_reconstruction}

\todo{Describe how the wind parameters (speed, direction) are estimated as part of the reconstruction.}

\subsection{Forcing maps features}
\label{sec:forcing_features}

\todo{Describe any priors/regularization imposed on the reconstructed maps (analogous to FMM section~2.4 use of external data, or CMM's parametric priors in Table~5).}

\section{Results}
\label{sec:results}

\todo{Present the results obtained from end-to-end simulations: reconstructed atmosphere maps, reconstructed sky maps after atmosphere removal, reconstructed wind, and residuals/noise characterization (see FMM section~3 and CMM section~7 for the analogous presentation).}

\section{Alternative method: successive map-making}
\label{sec:successive_mapmaking}

\todo{Present the alternative approach where the atmosphere and the sky are reconstructed in two successive map-making steps rather than jointly, and compare its performance to the joint approach of section~\ref{sec:algorithm}.}

\section{Conclusions}
\label{sec:conclusions}

\todo{Summarize the main results and their implications for QUBIC's B-mode measurement.}

\todo{Further developments: list planned improvements (e.g.\ more realistic atmosphere model, combination with spectral-imaging component separation from CMM, etc.).}

\acknowledgments
QUBIC is funded by the following agencies. France: ANR (Agence Nationale de la
Recherche) contract ANR-22-CE31-0016, DIM-ACAV (Domaine d’Intérêt Majeur-Astronomie et Conditions d’Apparition de la Vie), CNRS/IN2P3 (Centre national de la recherche scientifique/Institut national de physique nucléaire et de physique des particules), CNRS/INSU (Centre national de la recherche scientifique/Institut national et al de sciences de l’univers). Italy: CNR/PNRA (Consiglio Nazionale delle Ricerche/Programma Nazionale Ricerche in
Antartide) until 2016, INFN (Istituto Nazionale di Fisica Nucleare) since 2017.  Argentina: MINCyT (Ministerio de Ciencia, Tecnología e Innovación), CNEA (Comisión Nacional de Energía Atómica), CONICET (Consejo Nacional de Investigaciones Científicas y Técnicas).

\bibliographystyle{JHEP}
\bibliography{biblio}

\appendix

\section{\todo{Appendix title, e.g.\ Atmospheric correlations: Imager vs BI, if moved out of section~\ref{sec:atm_correlations}}}
\label{sec:appendixA}

\todo{Appendix content.}

\end{document}
